% GENERATED FILE. Edit canonical lessons, then run npm run generate:books.
% canonical-source-hash: fnv1a64:31e96db1322ec83a
% canonical-lessons: JA-C518-torimodosu, JA-C518-koriru, JA-C518-oshii, JA-C518-shikujiru, JA-C518-ochido

\chapter{Slips and Second Chances}
\label{ch:ja-slips-and-second-chances}
\hlchaptermodality{518}

\begin{chapteropening}
\textbf{By the end of this chapter:} \emph{I can say what would happen if things went wrong.}

Complete the last lesson of chapter 518: I can say what would happen if things went wrong.
\end{chapteropening}

\section[\texorpdfstring{\ja{とりもどす} (torimodosu)}{とりもどす (torimodosu)}]{\ja{とりもどす} (torimodosu) — to get back; to regain}

\label{lesson:JA-C518-torimodosu}



\begin{quote}
Before the new one: say the Japanese for to judge, then the Japanese for an option.
\end{quote}

\subsection*{You'll want to know: \ja{とりもどす}}
\textbf{\ja{とりもどす}} — \emph{torimodosu} — \textquotedblleft{}to get back; to regain\textquotedblright{}.

\begin{quote}
\textbf{\ja{もし} \ja{びょうきが} \ja{なおったら、かれは} \ja{げんきを} \ja{とりもどすでしょう。}} — \emph{moshi byōki ga naottara, kare wa genki o torimodosu deshō} — \textquotedblleft{}If his illness gets better, he will get his strength back.\textquotedblright{}
\end{quote}

Polite \textbf{\ja{とりもどします}}, \emph{torimodoshimasu}; a godan verb. Past (polite): \textbf{\ja{かれが} \ja{げんきを} \ja{とりもどしました。}}, \emph{kare ga genki o torimodoshimashita}.

\begin{cousinweb}
\textbf{\ja{とりもどす}} (\emph{torimodosu}) is a native Japanese word (wago), inherited rather than borrowed from Chinese, built from \emph{toru}, \textquotedblleft{}to take\textquotedblright{}, and \emph{modosu}, \textquotedblleft{}to put back\textquotedblright{}; its kanji was attached later as a kun'yomi (JMdict, KANJIDIC2).
\end{cousinweb}

\begin{grammarlens}[title={What you've built}]
One more word for saying what happens if.
\end{grammarlens}

\subsection*{Your turn}
\begin{itemize}
\raggedright
  \item \emph{Say it:} \emph{torimodosu}
  \item \emph{Say it:} \emph{torimodosu}, once more
  \item \emph{Say it:} \emph{moshi byōki ga naottara, kare wa genki o torimodosu deshō} — If his illness gets better, he will get his strength back.
  \item \emph{Recall:} say the Japanese for to judge, then the Japanese for an option, then say \emph{torimodosu} again
\end{itemize}

\begin{tcolorbox}[breakable,colback=teal!4,colframe=teal!35!black,title={Before you move on}]
What does \ja{とりもどす} mean? (\textquotedblleft{}To get back; to regain\textquotedblright{}.) Say it once more.
\end{tcolorbox}
\section[\texorpdfstring{\ja{こりる} (koriru)}{こりる (koriru)}]{\ja{こりる} (koriru) — to learn one's lesson}

\label{lesson:JA-C518-koriru}



\begin{quote}
Before the new one: say the Japanese for an option, then the Japanese for to get back.
\end{quote}

\subsection*{You'll want to know: \ja{こりる}}
\textbf{\ja{こりる}} — \emph{koriru} — \textquotedblleft{}to learn one's lesson\textquotedblright{}.

\begin{quote}
\textbf{\ja{いちど} \ja{けがを} \ja{したら、おとうとも} \ja{こりるでしょう。}} — \emph{ichido kega o shitara, otōto mo koriru deshō} — \textquotedblleft{}Once he gets hurt, my younger brother will learn his lesson too.\textquotedblright{}
\end{quote}

Polite \textbf{\ja{こります}}, \emph{korimasu}; an ichidan verb. Past (polite): \textbf{\ja{おとうとが} \ja{こりました。}}, \emph{otōto ga korimashita}.

\begin{cousinweb}
\textbf{\ja{こりる}} (\emph{koriru}) is a native Japanese word (wago), inherited rather than borrowed from Chinese; its kanji was attached later as a kun'yomi, and its older origin is not settled (JMdict, KANJIDIC2).
\end{cousinweb}

\begin{grammarlens}[title={What you've built}]
One more word for saying what happens if.
\end{grammarlens}

\subsection*{Your turn}
\begin{itemize}
\raggedright
  \item \emph{Say it:} \emph{koriru}
  \item \emph{Say it:} \emph{koriru}, once more
  \item \emph{Say it:} \emph{ichido kega o shitara, otōto mo koriru deshō} — Once he gets hurt, my younger brother will learn his lesson too.
  \item \emph{Recall:} say the Japanese for an option, then the Japanese for to get back, then say \emph{koriru} again
\end{itemize}

\begin{tcolorbox}[breakable,colback=teal!4,colframe=teal!35!black,title={Before you move on}]
What does \ja{こりる} mean? (\textquotedblleft{}To learn one's lesson\textquotedblright{}.) Say it once more.
\end{tcolorbox}
\section[\texorpdfstring{\ja{おしい} (oshii)}{おしい (oshii)}]{\ja{おしい} (oshii) — so close; regrettable}

\label{lesson:JA-C518-oshii}



\begin{quote}
Before the new one: say the Japanese for to get back, then the Japanese for to learn one's lesson.
\end{quote}

\subsection*{You'll want to know: \ja{おしい}}
\textbf{\ja{おしい}} — \emph{oshii} — \textquotedblleft{}so close; regrettable\textquotedblright{}.

\begin{quote}
\textbf{\ja{もし} \ja{いちてんで} \ja{まけたら、おしいでしょう。}} — \emph{moshi ichiten de maketara, oshii deshō} — \textquotedblleft{}If we lose by one point, it will be a real pity.\textquotedblright{}
\end{quote}

An i-adjective.

\begin{cousinweb}
\textbf{\ja{おしい}} (\emph{oshii}) is a native Japanese word (wago), inherited rather than borrowed from Chinese; its kanji was attached later as a kun'yomi, and its older origin is not settled (JMdict, KANJIDIC2).
\end{cousinweb}

\begin{grammarlens}[title={What you've built}]
One more word for saying what happens if.
\end{grammarlens}

\subsection*{Your turn}
\begin{itemize}
\raggedright
  \item \emph{Say it:} \emph{oshii}
  \item \emph{Say it:} \emph{oshii}, once more
  \item \emph{Say it:} \emph{moshi ichiten de maketara, oshii deshō} — If we lose by one point, it will be a real pity.
  \item \emph{Recall:} say the Japanese for to get back, then the Japanese for to learn one's lesson, then say \emph{oshii} again
\end{itemize}

\begin{tcolorbox}[breakable,colback=teal!4,colframe=teal!35!black,title={Before you move on}]
What does \ja{おしい} mean? (\textquotedblleft{}So close; regrettable\textquotedblright{}.) Say it once more.
\end{tcolorbox}
\section[\texorpdfstring{\ja{しくじる} (shikujiru)}{しくじる (shikujiru)}]{\ja{しくじる} (shikujiru) — to slip up; to fail (at something)}

\label{lesson:JA-C518-shikujiru}



\begin{quote}
Before the new one: say the Japanese for to learn one's lesson, then the Japanese for so close.
\end{quote}

\subsection*{You'll want to know: \ja{しくじる}}
\textbf{\ja{しくじる}} — \emph{shikujiru} — \textquotedblleft{}to slip up; to fail (at something)\textquotedblright{}.

\begin{quote}
\textbf{\ja{もし} \ja{しけんで} \ja{しくじったら、らいねん} \ja{もう} \ja{いちど} \ja{うけます。}} — \emph{moshi shiken de shikujittara, rainen mō ichido ukemasu} — \textquotedblleft{}If I slip up in the exam, I will take it again next year.\textquotedblright{}
\end{quote}

Polite \textbf{\ja{しくじります}}, \emph{shikujirimasu}; a godan verb. Past (polite): \textbf{\ja{おとうとが} \ja{しけんで} \ja{しくじりました。}}, \emph{otōto ga shiken de shikujirimashita}.

\begin{cousinweb}
\textbf{\ja{しくじる}} (\emph{shikujiru}) is usually counted a native Japanese word (wago), inherited rather than borrowed from Chinese: JMdict gives it no kanji spelling, and no dictionary we use records an older form, so its history is not settled (JMdict).
\end{cousinweb}

\begin{grammarlens}[title={What you've built}]
One more word for saying what happens if.
\end{grammarlens}

\subsection*{Your turn}
\begin{itemize}
\raggedright
  \item \emph{Say it:} \emph{shikujiru}
  \item \emph{Say it:} \emph{shikujiru}, once more
  \item \emph{Say it:} \emph{moshi shiken de shikujittara, rainen mō ichido ukemasu} — If I slip up in the exam, I will take it again next year.
  \item \emph{Recall:} say the Japanese for to learn one's lesson, then the Japanese for so close, then say \emph{shikujiru} again
\end{itemize}

\begin{tcolorbox}[breakable,colback=teal!4,colframe=teal!35!black,title={Before you move on}]
What does \ja{しくじる} mean? (\textquotedblleft{}To slip up; to fail (at something)\textquotedblright{}.) Say it once more.
\end{tcolorbox}
\section[\texorpdfstring{\ja{おちど} (ochido)}{おちど (ochido)}]{\ja{おちど} (ochido) — a fault; an oversight}

\label{lesson:JA-C518-ochido}



\begin{quote}
Before the new one: say the Japanese for so close, then the Japanese for to slip up.
\end{quote}

\subsection*{You'll want to know: \ja{おちど}}
\textbf{\ja{おちど}} — \emph{ochido} — \textquotedblleft{}a fault; an oversight\textquotedblright{}.

\begin{quote}
\textbf{\ja{もし} \ja{みせに} \ja{おちどが} \ja{あったら、あやまります。}} — \emph{moshi mise ni ochido ga attara, ayamarimasu} — \textquotedblleft{}If the shop is at fault, we will apologize.\textquotedblright{}
\end{quote}

\begin{cousinweb}
\textbf{\ja{おちど}} (\emph{ochido}) mixes the two layers, native and Sino-Japanese: native \emph{ochi}, \textquotedblleft{}falling\textquotedblright{}, from \emph{ochiru}, \textquotedblleft{}to fall\textquotedblright{}, and Sino-Japanese \emph{do}, \textquotedblleft{}degrees\textquotedblright{} (Middle Chinese \emph{duH}), borrowed from Chinese (JMdict, KANJIDIC2; Middle Chinese in Baxter's transcription).
\end{cousinweb}

\begin{grammarlens}[title={What you've built}]
One more word for saying what happens if.
\end{grammarlens}

\subsection*{Your turn}
\begin{itemize}
\raggedright
  \item \emph{Say it:} \emph{ochido}
  \item \emph{Say it:} \emph{ochido}, once more
  \item \emph{Say it:} \emph{moshi mise ni ochido ga attara, ayamarimasu} — If the shop is at fault, we will apologize.
  \item \emph{Recall:} say the Japanese for so close, then the Japanese for to slip up, then say \emph{ochido} again
\end{itemize}

\begin{tcolorbox}[breakable,colback=teal!4,colframe=teal!35!black,title={Before you move on}]
What does \ja{おちど} mean? (\textquotedblleft{}A fault; an oversight\textquotedblright{}.) Say it once more.
\end{tcolorbox}
